% !TEX root = ../../../main.tex
% Block card: detection.threshold -- "Threshold to spans" (Part I, the blocks chapter).
\cardfield{What it does}
Turns a curve of scores into stretches of time: every unbroken run where the score is above a
level you set becomes one span.

\cardfield{When to use it}
After any block that gives one score per sample, to get detections you can review and count: after
\blockref{detection.fluss} to get one span per change of behaviour, after
\blockref{detection.matrix_profile} to mark the stretches with no close look-alike, or after a
classifier's scores. It is the common last step of many detection chains.

\cardfield{What goes in, what comes out}
In goes a \term{scores}{scores} curve. Out come \term{span-set}{spans}, each scored by the highest
value inside it, written as detections.

\cardfield{Settings that matter}
\emph{threshold} is the level; only values strictly above it count. \emph{merge\_gap\_s} joins two
spans separated by a shorter gap than this, so an event broken by one dip below the level stays
one event. \emph{min\_duration\_s} then drops spans shorter than this, so a single sample that
grazes the level does not count. Both are off at 0.

\cardfield{Watch out for}
It only keeps what is \emph{above} the level. On a matrix profile, high means ``unlike anything
else'', so this block finds the oddities, not the repeats. The level is in the units of whatever
score comes in, and for matrix-profile scores it changes meaning with the window length: re-choose
it whenever the window changes. And when the scores belong to windows that start at each sample,
each span marks where those windows \emph{start}; the stretch of signal they cover runs on by up to
one window past the span's end.
